% !TEX TS-program = xelatex
% !TEX encoding = UTF-8 Unicode
% !Mode:: "TeX:UTF-8"

\documentclass{resume}
\usepackage{zh_CN-Adobefonts_external} % Simplified Chinese Support using external fonts (./fonts/zh_CN-Adobe/)
% \usepackage{NotoSansSC_external}
% \usepackage{NotoSerifCJKsc_external}
% \usepackage{zh_CN-Adobefonts_internal} % Simplified Chinese Support using system fonts
\usepackage{linespacing_fix} % disable extra space before next section
\usepackage{cite}

\begin{document}
\pagenumbering{gobble} % suppress displaying page number

\name{江紫阳}


\basicInfo{
  \email{jzysheepfoo@163.com} \textperiodcentered\ 
  \phone{(+86) 15851593862} \textperiodcentered\ 
  }

\section{教育背景}
\datedsubsection{\textbf{德克萨斯大学奥斯汀分校}, 奥斯汀}{2014年9月 -- 2016年5月}
\textit{工学硕士}\ 电子与计算机工程，绩点3.93/4.00
\datedsubsection{\textbf{南京理工大学}, 南京}{2009年9月 -- 2014年6月}
\textit{工学学士}\ 自动化，绩点90.27/100 排名3/175


\section{工作经历}
\datedsubsection{\textbf{Los Altos Hills Trading} 美国}{2023年5月 -- 2025年2月}
创始人
\begin{itemize}
  \item 一个交易信号订阅网站，和stripe集成支持支付
\end{itemize}

\datedsubsection{\textbf{Pure Storage} 山景城，加州}{2017年12月 -- 2022年11月}
研发工程师 2, 3, 4
\begin{itemize}
  \item 移动快照组和iscsi网络组成员
  \item 工作过快照到s3单元测试系统，s3客戶端，快照到nfs和s3垃圾处理代码重构，硬盘哈希等项目
  \item 研究了libcurl和boost asio源码，尤其是libcurl的域名系统缓存，连接池缓存，时钟树，boost asio的各种线和队列的实现
  \item 负责purevol收发项目，检查设计文档与测试方案，与不同组沟通，最终在flasharray大组下完成项目介绍
  \item 为移动快照组转移至欧洲做了一系列介绍
  \item 客戶端中断问题诊断的主要成员，参与了volume/vlan捆绑项目的前端swagger UI, python和中端java代码，帮助网络组别的成员的项目和代码审核
\end{itemize}

\datedsubsection{\textbf{Cadence Design Systems} 奥斯汀，德州}{2016年6月 -- 2017年12月}
软件工程师 2
\begin{itemize}
  \item 负责集成电路分块的内存预测函数以便更好散布在不同机器上做功耗的计算
  \item 参与了很多代码基础建设：libnetlist, 命名管理和功耗网络数据库
\end{itemize}

\datedsubsection{\textbf{The University of Texas at Austin} 奥斯汀，德州}{2015年8月 -- 2015年12月}
助教
\begin{itemize}
  \item EE380L 数据挖掘课的助教 
\end{itemize}

\datedsubsection{\textbf{Intersil} 奥斯汀，德州}{2015年5月 -- 2015年8月}
应用工程师实习
\begin{itemize}
  \item 用Python测试电源管理芯片
\end{itemize}


% Reference Test
%\datedsubsection{\textbf{Paper Title\cite{zaharia2012resilient}}}{May. 2015}
%An xxx optimized for xxx\cite{verma2015large}
%\begin{itemize}
%  \item main contribution
%\end{itemize}

\section{技能}
% increase linespacing [parsep=0.5ex]
\begin{itemize}[parsep=0.5ex]
  \item 编程语言:  C++, Java, Python
  \item 其他: AWS S3, Docker, Kubenetes, 分布式系统，paxos, raft, spanner, Cassandra, Memcached, Zookeeper
\end{itemize}

\section{获奖情况}
\datedline{江苏省三好学生}{2012年5月}
\datedline{国家奖学金}{2011年11月}

%% Reference
%\newpage
%\bibliographystyle{IEEETran}
%\bibliography{mycite}
\end{document}
